% Imported from topics/video-generation/seedance-2.5/long-video-spec.md
\chapter{长视频执行与因果缓存协议}\label{seedance-m24}
\moduleinfo{M24}{候选设计与工程选型；组合效果待验证}{\hyperref[src-long-context-tuning]{\texttt{long-context-tuning}}、\hyperref[src-longlive]{\texttt{longlive}}、\hyperref[src-longlive-code]{\texttt{longlive-code}}、\hyperref[src-ltx-2]{\texttt{ltx-2}}、\hyperref[src-ltx-original-code]{\texttt{ltx-original-code}}、\hyperref[src-ltx-original-model]{\texttt{ltx-original-model}}}

上一章明确了缓存成立的条件，这里把联合长片与因果续写分别写成候选执行过程。

长视频到这里分成两条计算图。一次联合生成时，未来镜头可反过来约束前面；逐块续写时，已交付历史却不能再依赖未来带噪状态。因此两者不只是请求长度不同，注意力、训练历史与缓存都要按这种依赖重新定义。

候选路线先训练联合长片，再派生块间因果版本。长上下文调优和流式方法提供可借鉴机制，但音视频联合迁移与训练尚未完成。下面先把需要满足的关系写清楚。

\section{共同的镜头状态是什么}\label{seedance-m24-ux5171ux540cux7684ux955cux5934ux72b6ux6001ux662fux4ec0ux4e48}

每个场景保存全局角色与环境描述，每个镜头保存局部事件、相机与时间范围；额外记录角色位置、服装、道具归属、已发生事件和待发生事件。

状态也有确定程度：用户指定、画面观察和模型推断分别记录证据时间与置信度。比如计划写钥匙在 A 手里，实际生成却已掉落，不能继续把计划当作发生事实；应重做结果或修正状态，再生成后续。

状态文本能帮助检索条件，却替代不了人物外观、动作和几何。新镜头仍要同时读取身份参考、相关关键帧和近期事件证据。若只把摘要越写越长，可能文字完整了，视觉条件仍然缺失。

\section{联合长片保留全局连接}\label{seedance-m24-ux8054ux5408ux957fux7247ux4fddux7559ux5168ux5c40ux8fdeux63a5}

第一版按8、16、30秒课程训练，视频与音频的token时间来自相同媒体时钟。每个样本使用全局条件和shot局部条件，新支路根据镜头与时间读取参考。

初始全段共享sigma，随后单列异步镜头任务。异步任务为每个镜头采样 \(\sigma_s\)，镜头覆盖内的视频和音频使用匹配的噪声状态；跨镜头边界token单列归属规则，不能同一token同时写两个sigma。

每个 token 的噪声状态也要真实传给模型。干净镜头为 0，活动镜头随步数变化；如果输入已异步，调制仍广播全局 sigma，两者就在描述不同状态。联合版本每步重算两条动态流，暂不把跨步历史 KV 当作精确值。

\section{因果版本如何分块}\label{seedance-m24-ux56e0ux679cux7248ux672cux5982ux4f55ux5206ux5757}

时间块目标约4秒，实际边界根据codec网格生成。视频与音频token各自保存位置网格区间，使用共同时间边界分配block ID；跨边界token延后归入能够完整容纳其区间终点的块，边界形成一个可报告的缓冲区。

时间网格还不是完整感受野。某个 token 的编码或预处理若读取未来，就要按真实依赖终点延后入块，或改用匹配的因果处理。只修改 block ID，没有消除编码阶段已经发生的信息泄漏。

因此按真实区间分块，比直接拿 token 下标除固定长度更清楚。训练推理用同一映射，RoPE 保留全局坐标；不能每次续写都将活动块当成新首帧，也不能把独立编码的局部首帧规则当作全局连续编码。

令query与key的块编号分别为 \(b_q,b_k\)：

\[
\operatorname{allow}(q,k)=\mathbf1[b_k\le b_q].
\]

同块内部双向，跨块只读过去，这条规则同时适用于视频 self、音频 self、视频读音频和音频读视频。若只改视频自身，未来仍可先进入声音，再从 AV 路径回来，因果约束并没有成立。

极简布尔关系见 \hyperref[code-block-causal-mask]{block\_causal\_mask.py}。真实长序列应使用块稀疏表示；该示例构造稠密布尔矩阵，只便于看清方向。

\section{原生前向接口还需要补什么}\label{seedance-m24-ux539fux751fux524dux5411ux63a5ux53e3ux8fd8ux9700ux8981ux8865ux4ec0ux4e48}

原Attention层支持mask参数，但原联合块的self与AV调用没有透传完整的候选因果mask。需要给 \codepath{TransformerArgs} 或额外执行上下文增加四类mask，并逐调用传入相应kernel。

其它支路也沿依赖检查：源条件塔、历史池化、字幕与对白能看到什么？用户预先给完整剧情或固定音轨，是已知任务条件；读取尚未生成的真实未来视频，则改变了预测问题。两者都叫“未来信息”会掩盖关键区别。

历史token噪声时间固定为0，当前块按采样sigma变化。跨模态调制使用各自token的噪声状态，避免把当前块的全局sigma广播进已缓存历史。

\section{缓存写入必须等到干净历史成立}\label{seedance-m24-ux7f13ux5b58ux5199ux5165ux5fc5ux987bux7b49ux5230ux5e72ux51c0ux5386ux53f2ux6210ux7acb}

第 \(k\) 块采样结束后，它成为下一块的干净条件。最后一个带噪步生成的隐藏状态仍属于那个带噪输入，不能直接写为永久历史缓存。

\begin{verbatim}
1. 用已有干净历史和当前噪声块完成当前块采样。
2. 固定生成结果、全局坐标与已确认状态，当前块sigma设为0。
3. 以同一因果执行图重新计算当前块的干净隐藏状态。
4. 写入各层流内与AV读取需要的K/V，提交历史版本。
5. 下一个块每步只计算活动块，并读取已提交历史。
\end{verbatim}

每层至少区分视频self、音频self、音频作为视频条件、视频作为音频条件四套K/V。它们有不同投影与调制，不能把一个名为 \codepath{history_kv} 的张量供所有模块共用。

历史只能读更早历史和固定条件，不读当前未来目标。这样同一份历史 K/V 对所有活动去噪步才保持不变，缓存的等价性有了计算图依据，而不只是实现里存下了一个张量。

\section{提示改变时怎样失效}\label{seedance-m24-ux63d0ux793aux6539ux53d8ux65f6ux600eux6837ux5931ux6548}

缓存键包括模型与适配器、历史内容、坐标、噪声策略、文本条件、参考、结构条件和状态版本。一个历史层若读取了新变化的条件，其K/V就需要刷新；不能仅按原始视频哈希命中。

首版可先保守处理：编辑历史或改变全局条件，就重建受影响前缀与后续。若希望新提示只作用未来，必须固定历史写入时所用条件，将新提示单独送给活动块，并在训练中使用同样规则。这是另一种明确的执行含义，不能靠缓存命中自动获得。

窗口还有一个容易漏掉的关系。旧 token 已被淘汰，保留的隐藏状态却可能早已读过它。因此按相同窗口策略逐步重放，和只拿当前窗口从头算一次，不一定相等。先用完整历史定义正确性，再独立比较有限窗口的收益和退化，才知道省掉了什么。

\section{自生成历史不能直接接任意真实未来}\label{seedance-m24-ux81eaux751fux6210ux5386ux53f2ux4e0dux80fdux76f4ux63a5ux63a5ux4efbux610fux771fux5b9eux672aux6765}

训练历史先从真实数据开始，再加入接近模型失真的受控扰动，最后才逐步使用学生自己生成的历史。每一步都在改变当前块所面对的条件分布，需要分别观察。

例如学生已生成钥匙掉落，原片后续却假设钥匙仍在手里，这时直接接原真实未来，监督就与条件冲突。可以让固定 teacher 在同一历史、参考和下一事件条件下重新标后续，再滤掉矛盾与严重错误；也可采用对齐的短 rollout 蒸馏。首版选前一项，因此 teacher 数据生成也要计入预算。

历史停止反传，当前块计算flow损失。逐步把自生成历史比例从30\%提高到70\%只是课程起点；若teacher重标质量不够，先维持真实/扰动历史，不把任意学生历史与真实未来拼成``可训练对''。

\section{解码与交付也有因果性}\label{seedance-m24-ux89e3ux7801ux4e0eux4ea4ux4ed8ux4e5fux6709ux56e0ux679cux6027}

再往输出端看，Transformer 能缓存，不等于 decoder 已经零前瞻。视频解码器可能读后续 latent，音频 vocoder 也有上下文和缓冲，所以“模型逐块生成”尚不能直接推出“帧可以立刻提交”。

首版先离线逐块生成、整段解码。若要在线输出，就另定义解码 halo 与前瞻，等需要的上下文到齐再提交；若要求完全无前瞻，则要相应因果 codec 及训练。KV-cache 本身没有改变解码器的未来依赖。

\section{长片验收按时间段展开}\label{seedance-m24-ux957fux7247ux9a8cux6536ux6309ux65f6ux95f4ux6bb5ux5c55ux5f00}

分别检查前段、中段和末段的身份、运动、道具状态、事件顺序与声画同步。联合三十秒、块式续写到更长时长、多镜头返回同场景，是三组不同任务。

缓存正确性先比较相同因果执行图的完整重算与缓存结果，再比较窗口和量化版本。真实实现可读 \href{https://github.com/NVlabs/LongLive/blob/31c70d9071cd7851d82939806dc9fac645fcabe5/LongLive1.0/pipeline/causal_inference.py}{LongLive因果推理}，LTX迁移入口为 \href{https://github.com/Lightricks/LTX-2/blob/bd92a5f408eb8c99f8c05e9fe2fcaaec8b56b638/packages/ltx-core/src/ltx_core/model/transformer/transformer.py}{联合块} 与 \href{https://github.com/Lightricks/LTX-2/blob/bd92a5f408eb8c99f8c05e9fe2fcaaec8b56b638/packages/ltx-core/src/ltx_core/model/transformer/attention.py}{Attention}。后者尚需上述因果化改动。
